% CV of Efstratios Tsoukanis -- compile with: pdflatex Resume.tex
\documentclass[11pt,letterpaper]{article}

\usepackage[T1]{fontenc}
\usepackage[utf8]{inputenc}
\usepackage{charter}
\usepackage[scaled=0.92]{helvet}
\usepackage{microtype}
\usepackage[margin=0.85in,top=0.7in,bottom=0.75in]{geometry}
\usepackage{xcolor}
\usepackage{titlesec}
\usepackage{enumitem}
\usepackage[hidelinks]{hyperref}

\definecolor{accent}{HTML}{2A7896}
\pagestyle{empty}
\setlength{\parindent}{0pt}

% Section headings: small caps in the accent color with a thin rule underneath
\titleformat{\section}{\large\scshape\color{accent}}{}{0pt}{}[\vspace{-0.6ex}{\color{accent!60}\titlerule[0.6pt]}]
\titlespacing*{\section}{0pt}{2.2ex}{1.2ex}

\setlist{nosep,leftmargin=1.2em,topsep=0.3ex}
\setlist[enumerate]{label=\arabic*.,leftmargin=1.8em,itemsep=0.6ex}

% \entry{bold left}{right-aligned date}{second line}
\newcommand{\entry}[3]{%
  \textbf{#1}\hfill{\small #2}\par
  \ifx&#3&\else #3\par\fi
  \vspace{0.8ex}}
\newcommand{\me}{\textbf{E.~Tsoukanis}}
\newcommand{\pubsub}[1]{\vspace{0.4ex}\textit{#1}\par\vspace{0.2ex}}

\begin{document}

{\centering
  {\fontsize{24}{28}\selectfont Efstratios Tsoukanis}\par\vspace{1.2ex}
  Visiting Scholar, Institute of Mathematical Sciences, Claremont Graduate University\par
  \vspace{0.4ex}
  \href{mailto:efstratios.tsoukanis@cgu.edu}{efstratios.tsoukanis@cgu.edu}
  \enspace$\cdot$\enspace +30 697 820 3648
  \enspace$\cdot$\enspace \href{https://etsoukan.github.io}{etsoukan.github.io}
  \enspace$\cdot$\enspace \href{https://scholar.google.com/citations?user=lJYEMRAAAAAJ&hl=en}{Google Scholar}\par
}

\section{Positions}
\entry{Visiting Scholar}{Current}{Institute of Mathematical Sciences, Claremont Graduate University, Claremont, CA}
\begin{itemize}
  \item Research with Distinguished Research Professor Hrushikesh Mhaskar on computational harmonic analysis and machine learning, supported by the US Office of Naval Research.
\end{itemize}
\vspace{0.8ex}
\entry{Research Assistant Professor}{2024--2025}{Institute of Mathematical Sciences, Claremont Graduate University, Claremont, CA}

\section{Education}
\entry{Ph.D.\ in Mathematics}{2024}{University of Maryland, College Park}
\begin{itemize}
  \item Dissertation: \textit{G-Invariant Representations using Coorbits}
  \item Advisor: Professor Radu Balan
\end{itemize}
\vspace{0.8ex}
\entry{M.S.\ in Mathematics}{2019}{University of Crete, Heraklion}
\begin{itemize}
  \item Thesis: \textit{Fourier multiplier and Fefferman counter-example}
  \item Advisor: Professor Themistoklis Mitsis
\end{itemize}
\vspace{0.8ex}
\entry{B.S.\ in Mathematics}{2016}{University of Crete, Heraklion}
\begin{itemize}
  \item Erasmus exchange at the University of Barcelona, 2014
\end{itemize}

\section{Research Interests}
Applied harmonic analysis and geometric machine learning: group-invariant representations
(embeddings of orbit spaces, injectivity and bi-Lipschitz properties), phase retrieval,
kernel methods and approximation theory.

\section{Publications}
\pubsub{Journal articles}
\begin{enumerate}
  \item R.~Balan and \me.
    G-invariant representations using coorbits: Bi-Lipschitz properties.
    \textit{Advances in Computational Mathematics} 52(5), 92, 2026.
    \href{https://arxiv.org/abs/2308.11784}{arXiv:2308.11784}
  \item N.~Dym, M.~Wellershoff, \me, D.~Levy, and R.~Balan.
    Quantitative bounds for sorting-based permutation-invariant embeddings.
    \textit{IEEE Transactions on Information Theory}, to appear, 2026.
    \href{https://arxiv.org/abs/2510.22186}{arXiv:2510.22186}
  \item H.~N.~Mhaskar, \me, and A.~D.~Jagtap.
    An approximation theory perspective on machine learning.
    \textit{Neural Networks}, 108841, 2026.
    \href{https://www.sciencedirect.com/science/article/pii/S0893608026003035}{ScienceDirect}
\end{enumerate}
\pubsub{Conference proceedings}
\begin{enumerate}[resume]
  \item H.~N.~Mhaskar, R.~O'Dowd, and \me.
    Active learning classification from a signal separation perspective.
    \textit{International Conference on Sampling Theory and Applications (SampTA)}, Vienna, Austria, 2025.
    \href{https://arxiv.org/abs/2502.16425}{arXiv:2502.16425}
  \item R.~Balan and \me.
    Relationships between the phase retrieval problem and permutation invariant embeddings.
    \textit{International Conference on Sampling Theory and Applications (SampTA)}, New Haven, CT, 2023.
    \href{https://arxiv.org/abs/2306.13111}{arXiv:2306.13111}
\end{enumerate}
\pubsub{Preprints}
\begin{enumerate}[resume]
  \item R.~Balan, \me, and M.~Wellershoff.
    Stability of sorting based embeddings.
    \href{https://arxiv.org/abs/2410.05446}{arXiv:2410.05446}, 2024.
  \item R.~Balan and \me.
    G-invariant representations using coorbits: Injectivity properties.
    \href{https://arxiv.org/abs/2310.16365}{arXiv:2310.16365}, 2023.
\end{enumerate}

\section{Talks}
\begin{itemize}[leftmargin=3.4em,labelwidth=2.8em,labelsep=0.6em,align=left,itemsep=0.6ex]
  \item[2025] \textit{Active learning classification from a signal separation perspective.}
    SampTA 2025, Vienna, Austria.
  \item[2023] \textit{Relationships between the phase retrieval problem and permutation invariant embeddings.}
    SampTA 2023, Yale University, New Haven, CT.
  \item[2023] \textit{Stability properties of coorbit embeddings.} CodEx Seminar.
  \item[2022] \textit{Noise stability of functions with low influences.}
    Computational Learning Theory and Fourier Analysis Summer School.
\end{itemize}

\section{Teaching}
\entry{Teaching Assistant, University of Maryland}{2018--2023}{Calculus II and III, Linear Algebra, Differential Equations}
\entry{Teaching Assistant, University of Crete}{2015--2016}{Calculus I, Analysis I}

\section{Service and Other Activities}
\begin{itemize}[itemsep=0.3ex]
  \item Referee for \textit{Foundations of Computational Mathematics} and \textit{Applied and Computational Harmonic Analysis}
  \item Research Interaction Team in Applied Harmonic Analysis: Fall 2021, Fall 2023
  \item Research Interaction Team in Deep Learning: Fall 2022, Spring 2023, Fall 2023
\end{itemize}

\section{Skills}
\textbf{Programming:} Python, MATLAB, \LaTeX \qquad \textbf{Languages:} Greek, English

\end{document}
